\section{Conclusiones}
La comparación muestra que el desempeño depende de la tarea: DQN alcanza el
umbral de retorno en persecución con menos interacciones que PPO, mientras
A2C/PPO obtienen mejores resultados de conducción después de imitación.
PBRS mejora el éxito frente al shaping ingenuo, cuya recompensa elevada no
se traduce en mejor cumplimiento de la tarea. BC y DAgger permiten aprovechar
demostraciones, pero su beneficio depende del experto y del ajuste posterior.

El mejor ataque jerárquico obtiene 130 goles en 200 episodios, equivalentes a
65,0\,\% (IC 95\,\%: 58,2--71,3\,\%), y supera en esta prueba a la jerarquía
inicial y a la política plana. La meta del 80\,\% permanece pendiente. Entrenar
definición dentro de la jerarquía ofrece una mejora relevante, sin aislar su
efecto del presupuesto adicional. Para consolidar la conclusión se requieren
réplicas con varias semillas, presupuestos comparables y condiciones de juego
más diversas.
